\subsection{库仑分支}
与希格斯分支相比，库仑分支更难描述。在较早的文献中，人们通常会看到这样的断言：「经典」库仑分支是 $T^*{}^LT/W$，其中 ${}^LT$ 是 $G$ 的极大环面 $T$ 的 Langlands 对偶。然而这并不是真正的答案，因为存在非平凡的「量子修正」。在某些特殊情形下，真正的库仑分支可以由其他方法确定：
\begin{itemize}
	\item Hanany--Witten \cite{hananyTypeIIB1997} 的工作及其推广，利用弦对偶确定了箭图规范理论 E1. 与 E2. 的库仑分支。
	\item 在 \cite{telemanGaugeTheory2014} 中隐含了这样一个结论：规范群为 $G$ 的纯规范理论（即 $N=0$）的库仑分支，是 $G^{\vee}$ 的（李代数）正则中心化子簇（regular centralizer variety），该文献中称之为 Bezrukavnikov--Finkelberg--Mirkovi\'c 空间 $\operatorname{BFM}(G^{\vee})$。这个空间被等同于 $G^{\vee}$ 的仿射 Grassmann 流形的等变同调的谱，见 \cite{bezrukavnikovEquivariantHomology2005}。
\end{itemize}


Braverman、Finkelberg 与 Nakajima 的工作 \cite{BFN} 基于仿射 Grassmann 流形的几何，给出了当 $X=T^*N$ 时库仑分支的一个显式数学定义。让我们试着粗略解释这个构造的来源，它看起来可能相当令人生畏。

按照描述（\ref{fns:A}），我们应当考虑从圆柱面到 $N/G$ 的映射：它们在圆柱的上下两个圆盘上全纯，并沿侧面为常值。我们将定义圆柱的参数 $\delta$ 收缩为无穷小，只考虑上下两个平面 $z=\pm \delta$ 中原点处的 Taylor 展开。若令 $t = x + iy$，则每个全纯映射 $D\to N/G$ 对应一个 Taylor 级数 $n(t)\in N[[t]]$；两个映射给出同一个映射，当且仅当它们处于 $G$ 值 Taylor 级数 $G[[t]]$ 的同一轨道中。

最后，映射沿圆柱的侧面应为常值。映射到商空间意味着：若两个对象处于圆上某个群值函数的同一轨道中，则它们相等。由于这只是在圆柱的侧面上，比较上下两侧的群值函数可能在原点处有极点。因此，若 $n_\pm(t)$ 是 $(0,0,\pm \delta)$ 处的 Taylor 展开，则必须存在某个 $G$ 值 Laurent 级数 $g_{\pm}(t)\in G((t))$，使得 $g_{\pm} n_{-}(t)=n_+(t)$。

\begin{definition}
    $(G,N)$ 的{\bf BFN 空间}（BFN space）是集合 $\{(n_+,n_-,g_{\pm}) \mid g_{\pm} n_{-}(t)=n_+(t)\}$ 在 $G[[t]]\times G[[t]]$ 作用下的商，作用为 \[(h_+,h_-)\cdot (n_+,n_-,g_{pm})=(h_+n_+,h_-n_-,h_+g_{\pm}h_-^{-1}) \]
\end{definition}
\begin{yizhu}[下标笔误]
此处照录原文公式。原文在等号左边写作 $g_{pm}$，而同一公式的其余各处以及上一段均作 $g_{\pm}$，且 $g_{\pm}$ 正是上一段引入的记号；故 $g_{pm}$ 应为 $g_{\pm}$ 之误。
\end{yizhu}
根据（\ref{fns:A}），我们应当考虑这个空间上的局部常值函数。这里必须小心：事实上，我们需要考虑这个商的 Borel--Moore 同调（须非常仔细地定义）。也就是说：
\begin{definition}\label{def:Coulomb}
    代数 $\alg_{\operatorname{A}}(T^*N,G)$ 是 BFN 空间的 Borel--Moore 同调。
\end{definition}
\begin{yizhu}[BFN 空间]
BFN 空间把（收缩后的）圆柱面上的场论数据约化为仿射 Grassmann 流形上的三元组 $(n_+,n_-,g_{\pm})$，再对 $G[[t]]\times G[[t]]$ 取商。由于这个作用不自由，商上的 Borel--Moore 同调不能按字面理解，须借助等变拓扑来定义；精确的处理见前引 Braverman--Finkelberg--Nakajima 的工作。
\end{yizhu}
虽然 $G[[t]]\times G[[t]]$ 的作用并不自由，我们仍可以用等变拓扑来解释这个商的同调；感兴趣的读者应参阅 \cite{BFN} 中更精确的讨论。我们照样能用有限维拓扑中的通常方法来计算，并在许多情形下确定库仑分支。

若 $X$ 是 $G$ 的辛表示，且无法写成 $X=T^*N$ 的形式（$N$ 为不变子空间），则这个定义必须修改，而且库仑分支的存在似乎有障碍。这一进路最近由 Braverman--Dhillon--Finkelberg--Raskin--Travkin \cite{bravermanCoulombBranches2022} 与 Teleman \cite{telemanCoulombBranches2022} 的工作发展出来。

库仑分支自带一个 $\C^{\times}$ 作用，它由 Borel--Moore 同调上的同调分级诱导。不幸的是，其中可能出现负次数的算子。我们称这样的理论为{\bf 坏}（bad）。大多数其他情形称为{\bf 好}（good）（尽管还有第三种可能，读者大概能猜到它叫「丑」（ugly），但我们将忽略这种情形）。

一个已经引起数学家兴趣的坏例子，是纯规范理论中 $N=0$ 的情形。此时希格斯分支相当退化，但 $\mathfrak{M}_{\operatorname{A}}=\operatorname{BFM}(G)$ 是有理 Toda 链（rational Toda lattice）的相空间，即限制在 Kostant 切片上的万有（李代数）中心化子 \cite{bezrukavnikovEquivariantHomology2005}；关于这个簇及其相关文献的更详细讨论，见 \cite[\S 5.1]{telemanGaugeTheory2014}。我们所知道的其他有趣情形，都属于前面讨论过的两类：

\subsubsection{阿贝尔/超环面规范理论}

若 $G$ 是阿贝尔群（abelian group）且忠实地作用于 $N$，则库仑分支将与另一个好的阿贝尔理论的希格斯分支重合。若作用不忠实，理论将是丑的。

回忆一下，我们可以假定 $G\subset D$ 是对角矩阵群 $D$ 的子群。$F=D/G$ 的 Langlands 对偶群 ${}^LF$ 可实现为 $D$ 的连通子群，其李代数是 $\mathfrak{g}\subset \mathfrak{d}=\C^n$ 的正交补；它也同构于 $(\C^{\times})^k$，但同构 $F\cong {}^LF$ 并非典范。

于是我们得到两个环面的短正合列，它们互为 Langlands 对偶：
\begin{align*}
	1\to G\to & D \to F \to 1\\
	1\to {}^LG\leftarrow & D\leftarrow {}^L F \leftarrow 1.
\end{align*}
相应的规范理论之间有如下关系：
\begin{theorem}
  我们有同构
\[ \mathfrak{M}_{\operatorname{A/B}}(T^*\C^n,G)\cong \mathfrak{M}_{\operatorname{B/A}}(T^*\C^n,{}^LF). \]
\end{theorem}
这个同构在库仑分支的精确数学定义出现之前就已被广泛预期。它似乎最早是由 Dimofte 与 Hilburn 用上述数学定义验证的；其物理讨论见 \cite[\S 3.3]{bullimoreCoulombBranch2017}，而优雅的数学证明见 \cite[\S 4(vii)]{BFN}。

虽然完整描述这个同构超出了本文的范围，我们还是想略述其风味。首先，两边都带有 $F$ 的作用；等价地，$\alg_{\operatorname{A/B}}$ 上带有按 $F$ 的权格给出的分级，或按 ${}^LF$ 的余权格给出的分级：
\begin{enumerate}
	\item 在 $\mathfrak{M}_{\operatorname{B}}(T^*\C^n,G)$ 上，它由 $T^*\C^n$ 上的 $D$ 作用诱导。
	\item $\alg_{\operatorname{A}}(T^*\C^n,{}^LF)$ 上所需的分级，由 ${}^LF$ 的仿射 Grassmann 流形的各连通分支到余权格的双射诱导。
\end{enumerate}
其次，在两个代数中，次数为 $0$ 的元素都构成 $\Sym(\mathfrak{f})$ 的一份拷贝：
\begin{enumerate}
	\item 在 $\mathfrak{M}_{\operatorname{B}}(T^*\C^n,G)$ 上，它们是经由矩映射 $\mathfrak{M}_{\operatorname{B}}(G,T^*\C^n)\to \mathfrak{f}^*$ 分解的多项式函数。
	\item 在 $\alg_{\operatorname{A}}(T^*\C^n,{}^LF)$ 中，次数为 $0$ 的元素是点上的 Borel--Moore 同调模去 ${}^LF[[t]]$ 的作用，这等于 $\Sym(\mathfrak{f})\cong H^*_{{}^LF}(\operatorname{pt})$。
\end{enumerate}
代数之间的同构已几乎由希格斯分支与库仑分支这两个方面的匹配所决定。


\subsubsection{箭图规范理论}
\label{sec:quiver-gauge-theories}

我们前面讨论过的另一类理论——也是我们目前理解得最丰富的——是箭图规范理论。由此得到的库仑分支正是 Nakajima 箭图簇的镜像。那些在上述意义下是好的理论，对应于一对数据：Dynkin 图由该箭图给出的 Kac--Moody 李代数的一个最高权 $\lambda$，以及一个支配权 $\mu$。它们由以下条件刻画：
\begin{enumerate}
    \item 最高权向量关于节点 $i$ 所对应的 $\mathfrak{sl}_2$ 的权为 $w_i$（即 $\alpha_i^{\vee}(\lambda)=w_i$）。
    \item 对单根 $\alpha_i$，有 $\lambda-\mu=\sum v_i\alpha_i.$
\end{enumerate}
当 $\Gamma$ 是 ADE Dynkin 图时，这些库仑分支可以用 Dynkin 图为 $\Gamma$ 的有限维群 $G_{\Gamma}$ 的仿射 Grassmann 流形来解释。
我们可以把 $\lambda$ 与 $\mu$ 解释为 $G_{\Gamma}$ 的余权，从而考虑轨道闭包 $\overline{\mathsf{Gr}}^{\lambda} =\overline{\GO t^{\lambda}\GO/\GO}\subset \GK/\GO$。
\begin{proposition}[\mbox{\cite[Th. 3.10]{BFNplus}}]
	ADE 箭图规范理论的库仑分支 $\mathfrak{M}_{\operatorname{A}}(T^*N,G)$，同构于 $\overline{\mathsf{Gr}}^{\lambda}$ 内 $\overline{\mathsf{Gr}}^{\mu}$ 的一般点处的横截切片。
\end{proposition}

通过引入仿射 Grassmann 流形中的某些广义切片，这个结果可以推广到非支配权 \cite{BFNplus}。由于这些仿射 Grassmann 流形切片对大多数读者来说并不熟悉，让我们讨论几个例子：



\begin{ex}
箭图 E1. 与 E2. 都满足 $\mathfrak{M}_{\operatorname{A}}\cong \mathfrak{M}_{\operatorname{B}}$！这是巧合，通常不会发生。
\end{ex}
\begin{ex}
    对箭图 \tikz[baseline=-2pt]{
		\node[draw, thick, circle, inner sep=2pt,fill=white] (d) at (0,0) {$1$};
	\node[draw, thick,inner sep=4pt] (s) at (1,0){$n$};
		\draw[thick] (d) -- (s) ;
}，其中 $\mathfrak{M}_{\operatorname{B}}\cong \rkone$，库仑分支 $\mathfrak{M}_{\operatorname{A}}$ 是仿射簇 $\C^2/\mathbb{Z}_{n}$，循环群 $\mathbb{Z}_{n}$ 通过矩阵 $\operatorname{diag}(e^{2\pi ik/n},e^{-2\pi ik/n})$ 作用。注意在此情形规范群是阿贝尔的，因此我们可以使用超环面描述（hypertoric description）。
\end{ex}
\begin{ex} 在下面的箭图中，希格斯分支与库仑分支与上一例正好互换：$\mathfrak{M}_{\operatorname{B}}\cong \C^2/\mathbb{Z}_{n}$，$ \mathfrak{M}_{\operatorname{A}}\cong \rkone$。同样，这一情形是超环面的，并与上一例对偶。
\begin{equation*}
\tikz{
	\node[draw, thick,inner sep=12pt]  (a) at (-4.5,0) {$1$};
	\node[draw, thick, circle, inner sep=6pt,fill=white] (b) at (-3,0) {$1$};	
		\node[draw, thick, circle, inner sep=6pt,fill=white] (d) at (0,0) {$1$};
	\node[draw, thick,inner sep=12pt] (s) at (1.5,0){$1$};
	\node[inner sep=12pt,fill=white] (c) at (-1.5,0){$\cdots$};
		\draw[thick] (a) -- (b) ;
		\draw[thick] (b) -- (c) ;
		\draw[thick] (c) -- (d) ;
		\draw[thick] (d) -- (s) ;
}
\end{equation*}
\end{ex}

这些都是一个远为一般的结果的特例。对于 $\Gamma$ 为仿射 A 型 Dynkin 图（即单个循环）的好箭图理论，包括 Jordan 箭图（单个自环），库仑分支也是某个仿射 A 型 Dynkin 图理论的希格斯分支，只是规模可能不同，这已在 \cite{nakajimaCherkisBow2017} 中证明。这个对应的组合学稍显复杂，但与已知的秩--水平对偶（rank-level duality）的组合学相符。这暗示相应的理论互为镜像。

A 型（即线性）箭图的情形是循环情形的特例（令某个 $v_i=0$ 就「破坏」了循环），此时镜像仍是一个线性箭图的箭图规范理论。例如，$T^*\operatorname{Gr}(w,v)$，即 $\C^w$ 中 $v$ 维平面构成的 Grassmann 流形的余切丛，是单节点箭图规范理论 \tikz[baseline=-2pt]{	
		\node[draw, thick, circle, inner sep=2pt,fill=white] (d) at (0,0) {$v$};
	\node[draw, thick,inner sep=4pt] (s) at (1,0){$w$};
		\draw[thick] (d) -- (s) ;
} 的消解希格斯分支；当 $2v\leq w$ 时，这个理论是好的。若 $2v<w$，则镜像理论由
\begin{equation*}
\tikz[scale=.9]{
	\node[draw, thick, inner sep=10pt] (bb) at (-6,1.5) {$1$};
	\node[draw, thick, circle, inner sep=6pt,fill=white] (a) at (-6,-3) {$1$};
	\node[inner sep=12pt,fill=white] (ab) at (-6,-1.5){$\vdots$};
	\node[draw, thick, circle, inner sep=6pt,fill=white] (b) at (-6,0) {$v$};	
	\node[draw, thick, inner sep=10pt] (dd) at (0,1.5) {$1$};
		\node[draw, thick, circle, inner sep=6pt,fill=white] (d) at (0,0) {$v$};
	\node[draw, thick, circle, inner sep=6pt,fill=white] (e) at (0,-3) {$1$};
	\node[inner sep=12pt,fill=white] (de) at (0,-1.5){$\vdots$};
	\node[inner sep=12pt,fill=white] (c) at (-3,0){$\cdots$};
		\draw[thick] (a) -- (ab) --(b) --(bb);
		\draw[thick] (b) -- (c) ;
		\draw[thick] (c) -- (d) ;
		\draw[thick] (e) -- (de) --(d) --(dd);
				\node[draw, thick, circle, inner sep=6pt,fill=white] (d) at (-1.5,0) {$v$};
						\node[draw, thick, circle, inner sep=6pt,fill=white] (d) at (-4.5,0) {$v$};
}
\end{equation*}
给出，其水平线上有 $w-2v+1$ 个标号为 $v$ 的节点；而若 $2v=w$，则镜像理论为
\begin{equation*}
\tikz[scale=.9]{
	\node[draw, thick, inner sep=10pt] (cc) at (-3,1.5) {$2$};
	\node[draw, thick, circle, inner sep=6pt,fill=white] (a) at (-6,-3.5) {$1$};
	\node[inner sep=12pt,fill=white] (ab) at (-6,-2){$\vdots$};
	\node[draw, thick, circle, inner sep=2pt,fill=white] (b) at (-6,-.5) {$v-1$};	
		\node[draw, thick, circle, inner sep=6pt,fill=white] (c) at (-3,0) {$v$};
		\node[draw, thick, circle, inner sep=2pt,fill=white] (d) at (0,-.5) {$v-1$};
	\node[draw, thick, circle, inner sep=6pt,fill=white] (e) at (0,-3.5) {$1$};
	\node[inner sep=12pt,fill=white] (de) at (0,-2){$\vdots$};
		\draw[thick] (a) -- (ab) --(b);
		\draw[thick] (b) -- (c) --(cc);
		\draw[thick] (c) -- (d) ;
		\draw[thick] (e) -- (de) --(d) ;
}
\end{equation*}

在此情形下，对偶的组合学可以这样更方便地描述：把希格斯分支与库仑分支都等同于横截切片 $\mathfrak{X}^{\nu}_{\eta}$——对 $n$ 的分拆 $\eta,\nu$，它是 Jordan 型为 $\eta$ 的幂零矩阵轨道在 Jordan 型为 $\nu$ 的轨道闭包中的横截切片；这总是可以做到的。此时同一理论的库仑分支是切片 $\mathfrak{X}^{\eta^t}_{\nu^t}$。在上面的 Grassmann 流形情形中，理论 \tikz[baseline=-2pt]{	
		\node[draw, thick, circle, inner sep=2pt,fill=white] (d) at (0,0) {$v$};
	\node[draw, thick,inner sep=4pt] (s) at (1,0){$w$};
		\draw[thick] (d) -- (s) ;
} 的希格斯分支是 Jordan 型为 $(2^v,1^{w-2v})$ 的幂零矩阵轨道的闭包（即 $T^*\operatorname{Gr}(w,v)$ 在到 $\mathfrak{gl}_{w}^*$ 的矩映射下的像），而库仑分支则是整个幂零锥中到 Jordan 型为 $(w-v,v)$ 的轨道的切片，通常称为{\bf 两行 Slodowy 切片}（two row Slodowy slice）；它在表示论与纽结同调中已多次出现。
% [原文 p15–p19]
